\begin{pastexam}[label={px:ccu-109-swa-04}]{中正資工所甲組 109 學年度 第 9 題}
（計算）

(4\%) Find the numbers of different binary trees with 2, 3, 4, and 5 nodes, respectively.
\end{pastexam}

\begin{pastsoln}
答案：$2$、$5$、$14$、$42$。

「不同的二元樹」指形狀不同（節點沒有標號），左子樹和右子樹要分開算，只有左孩子和只有右孩子是
兩棵不同的樹。 令 $b_n$ 是 $n$ 個節點的二元樹個數，$b_0 = 1$（空樹）。

\textbf{遞迴式}。 $n \ge 1$ 個節點的二元樹，根以外還有 $n - 1$ 個節點，左子樹 $k$ 個、右子樹
$n - 1 - k$ 個，$k = 0, 1, \ldots, n - 1$。 左右子樹各自是任意的二元樹，而且不同的 $k$ 或不同的子樹
組出不同的樹，所以
\[
b_n = \sum_{k=0}^{n-1} b_k\, b_{n-1-k}.
\]
逐一算：
\begin{align*}
b_1 &= b_0 b_0 = 1,\\
b_2 &= b_0 b_1 + b_1 b_0 = 1 + 1 = 2,\\
b_3 &= b_0 b_2 + b_1 b_1 + b_2 b_0 = 2 + 1 + 2 = 5,\\
b_4 &= b_0 b_3 + b_1 b_2 + b_2 b_1 + b_3 b_0 = 5 + 2 + 2 + 5 = 14,\\
b_5 &= b_0 b_4 + b_1 b_3 + b_2 b_2 + b_3 b_1 + b_4 b_0 = 14 + 5 + 4 + 5 + 14 = 42.
\end{align*}

\textbf{公式}。 這個遞迴式的解是卡特蘭數 $b_n = \dfrac{1}{n+1}\dbinom{2n}{n}$（資結書定理~7.6），
可以用來核對：
\[
\frac{1}{3}\binom{4}{2} = \frac{6}{3} = 2,\quad
\frac{1}{4}\binom{6}{3} = \frac{20}{4} = 5,\quad
\frac{1}{5}\binom{8}{4} = \frac{70}{5} = 14,\quad
\frac{1}{6}\binom{10}{5} = \frac{252}{6} = 42.
\]

$n = 3$ 的五棵可以直接畫出來檢查：根的左右子樹大小是 $(2, 0)$ 有兩棵（左子樹本身往左或往右）、
$(1, 1)$ 一棵、$(0, 2)$ 兩棵，共 $5$ 棵。
\end{pastsoln}
